\documentclass[11pt]{relatorio_tcc_ads_ifba}

% --- Pacotes e Configurações ---
% (Seus pacotes, como hyperref, graphicx, etc., devem estar aqui ou na classe)

\begin{document}

% --- Informações Pré-Textuais ---
\aluno{Gabriel Vinicius Souza de Souza}
\titulo{Design e Implementação de um Pipeline de Processamento de Mídia Genérico Baseado em Arquitetura Modular e Orientada a Eventos}
\date{Agosto de 2024}
\orientador{Prof. Dr. Jowaner Araújo}
\coorientador{}

% Gera a capa, folha de rosto, etc.
\pretextual

% ---------------------------------------------------------------------
% CORPO DO DOCUMENTO
% ---------------------------------------------------------------------

% Cada arquivo de capítulo é responsável por seus próprios `\input` de subseções.

\chapter{Introdução}
\label{sec:introducao}

Na era digital contemporânea, a produção e o consumo de conteúdo de mídia, especialmente vídeos e áudios, atingiram uma escala sem precedentes. Plataformas de educação, entretenimento e comunicação corporativa dependem cada vez mais da capacidade de extrair informações valiosas e insights de grandes volumes de dados não estruturados. Nesse contexto, a automação de tarefas como transcrição e sumarização de conteúdo tornou-se não apenas uma conveniência, mas uma necessidade estratégica. No entanto, a construção de sistemas capazes de realizar essas tarefas de forma eficiente, escalável e, acima de tudo, flexível, apresenta desafios arquiteturais significativos.

Sistemas de processamento de mídia tradicionais frequentemente adotam arquiteturas monolíticas e fortemente acopladas, onde o fluxo de processamento é rígido e focado em um único tipo de mídia, geralmente o vídeo. Essa abordagem resulta em sistemas frágeis e de difícil manutenção, onde a adição de um novo tipo de mídia (como um arquivo de áudio ou um documento PDF) ou a integração de um novo serviço de IA (como um provedor de transcrição alternativo) exige refatorações complexas e arriscadas. A falta de separação de responsabilidades e o acoplamento direto entre os componentes do pipeline limitam a escalabilidade e a resiliência do sistema como um todo.

Este trabalho aborda esses desafios através do design e da implementação de uma plataforma de processamento de mídia genérica, fundamentada em uma arquitetura modular, orientada a eventos e baseada em padrões de design de software modernos. A solução proposta abstrai o conceito de "vídeo" para uma entidade de "Mídia" mais genérica, permitindo que o sistema orquestre pipelines de processamento dinâmicos e desacoplados para diferentes formatos de entrada. A comunicação entre os módulos especializados (e.g., extração de áudio, transcrição, sumarização) é realizada exclusivamente através de um barramento de eventos, garantindo alta coesão e baixo acoplamento.

A relevância deste projeto reside na aplicação prática e na síntese de um conjunto de padrões de arquitetura avançados, como Command Query Responsibility Segregation (CQRS), o Padrão State para gerenciamento de ciclo de vida, e o Padrão Strategy para a seleção dinâmica de provedores de serviço. O resultado é uma plataforma que não apenas cumpre os requisitos funcionais de transcrição e sumarização, mas o faz de uma maneira que é inerentemente extensível, resiliente e manutenível, servindo como um estudo de caso para a construção de sistemas distribuídos modernos.

Este documento está estruturado da seguinte forma: o Capítulo \ref{sec:objetivos} detalha os objetivos geral e específicos do projeto. O Capítulo \ref{sec:definicoes} apresenta o problema, a proposta de solução e as tecnologias utilizadas. O Capítulo \ref{sec:requisitos} descreve os requisitos funcionais e não funcionais. O Capítulo \ref{sec:design} aprofunda-se no design do sistema, detalhando a arquitetura de software e o modelo de banco de dados. Os capítulos subsequentes abordam a qualidade, a implantação e o manual do usuário, concluindo com as considerações finais e os resultados alcançados.
 % Capítulo de Introdução

\section{Objetivos}
    O objetivo geral do projeto é desenvolver uma plataforma automática de resumos de vídeos capaz de extrair conteúdo relevante de vídeos, gerando transcrições textuais e resumos concisos através de técnicas de processamento de áudio, reconhecimento de fala e sumarização de texto. A solução visa facilitar o acesso a conteúdos audiovisuais por meio de representações textuais estruturadas, auxiliando estudantes, profissionais e pesquisadores no consumo rápido de informações.
    
    Os objetivos específicos incluem:
    
        
- Implementar funcionalidades de upload e armazenamento seguro de vídeos;
        
- Realizar transcrição automática de áudio extraído dos vídeos;
        
- Aplicar técnicas de inteligência artificial para geração de resumos textuais;
        
- Garantir autenticação segura de usuários e notificações sobre o status dos processamentos;
        
- Fornecer interface intuitiva para visualização de resultados e gerenciamento de vídeos. % Capítulo de Objetivos

\section{Definições}
    Este documento descreve o desenvolvimento de uma plataforma web responsável pela automatização de resumos de vídeos. O sistema utiliza tecnologias modernas de reconhecimento de fala (ASR - Automatic Speech Recognition) e sumarização de texto, combinadas com uma arquitetura modular e escalável. O termo “plataforma” refere-se ao conjunto de componentes lógicos e físicos envolvidos na coleta, processamento e apresentação de dados relacionados aos vídeos submetidos pelos usuários.
 % Capítulo de Definições e Conceitos

\section{Requisitos}
\label{sec:requisitos}

Nesta seção, são detalhados os requisitos funcionais e não funcionais que guiaram o desenvolvimento da plataforma. Os requisitos funcionais descrevem o que o sistema deve fazer da perspectiva do usuário, enquanto os requisitos não funcionais definem os atributos de qualidade e as restrições técnicas sob as quais o sistema deve operar, justificando as principais decisões arquiteturais do projeto.
 % Capítulo de Requisitos

\section{Design e Arquitetura do Sistema}
\label{sec:design}

O design e a arquitetura de um sistema de software são os alicerces sobre os quais sua funcionalidade, manutenibilidade e escalabilidade são construídas. Para um sistema distribuído e assíncrono como a plataforma de processamento de mídias, a escolha de uma arquitetura robusta não é apenas uma boa prática, mas um requisito fundamental para o sucesso do projeto. Decisões tomadas nesta fase impactam diretamente a capacidade do sistema de evoluir, de se adaptar a novas tecnologias e de resistir a falhas.

Esta seção detalha as decisões de design e a arquitetura de alto nível que foram implementadas. Começando com uma visão arquitetural abrangente, que descreve os padrões de design e os princípios que guiaram o desenvolvimento, passando pela modelagem visual do sistema com diagramas UML, e culminando no design do esquema do banco de dados que sustenta a persistência dos dados. O objetivo é fornecer uma compreensão clara e profunda de como o sistema foi estruturado para atender aos seus requisitos funcionais e, crucialmente, aos seus atributos de qualidade não funcionais, como extensibilidade, resiliência e observabilidade.
 % Capítulo de Design e Arquitetura

\section{Qualidade de Software e Testes}
\label{sec:qualidade}

A garantia de qualidade é um pilar fundamental na engenharia de software, especialmente em sistemas distribuídos e assíncronos como o proposto neste trabalho. A complexidade inerente a uma arquitetura orientada a eventos, com múltiplos componentes desacoplados, exige uma abordagem disciplinada e multifacetada para assegurar a corretude, a robustez e a confiabilidade do sistema.

Neste contexto, a estratégia de qualidade adotada não se limitou apenas à verificação funcional do software, mas foi integrada ao próprio processo de design e desenvolvimento. Foram empregadas práticas de análise estática de código, como 	extit{linting} e verificação de tipos, para garantir a consistência e a legibilidade. Além disso, uma estratégia de testes automatizados em múltiplas camadas foi implementada para validar o sistema em diferentes níveis de abstração, desde a lógica de negócio de componentes individuais até o fluxo de processamento completo da aplicação. A subseção a seguir detalha o plano de testes executado, que constitui a base para a validação da qualidade e da corretude da plataforma desenvolvida.
 % Capítulo de Qualidade e Testes

\section{Implantação}
\label{sec:implantacao}

A implantação de um sistema de software em um ambiente de produção é a etapa final que torna a aplicação acessível e utilizável. Para um sistema distribuído como a plataforma de processamento de mídias, a estratégia de implantação deve garantir não apenas a disponibilidade, mas também a escalabilidade, a segurança e a consistência do ambiente.

Esta seção detalha o plano de implantação do projeto, que é baseado em uma abordagem moderna de containerização com Docker e orquestração com Docker Compose. A estratégia adotada visa criar um ambiente de produção portátil, replicável e otimizado, separando claramente as configurações de desenvolvimento das de produção para garantir estabilidade e performance. A subseção a seguir descreve a arquitetura de implantação, incluindo a configuração dos contêineres, a rede e o uso de servidores de aplicação robustos.
 % Capítulo de Implantação

\section{Manual do Usuário}
\textcolor{blue}{
    O manual do usuário será fornecido em formato digital, contendo instruções passo a passo sobre:
    
        
- Como criar uma conta e fazer login
        
- Como fazer upload de vídeos
        
- Como acompanhar o progresso do processamento
        
- Como visualizar transcrições e resumos
        
- Como gerenciar seus vídeos salvos
    

    O sistema possuirá telas intuitivas e feedbacks claros em cada etapa do processo. Em versões futuras, pretende-se incluir tutoriais interativos e suporte via chatbot.
} \\
\textcolor{red}{[Obrigatório]} % Capítulo do Manual do Usuário (API)

\chapter{Conclusão}
\label{sec:conclusao}

Este trabalho apresentou o design e a implementação de uma plataforma de processamento de mídia genérica, cuja concepção foi guiada pela necessidade de superar as limitações de sistemas tradicionais, como o alto acoplamento e a baixa extensibilidade. O problema foi abordado em duas frentes: a do usuário final, que enfrenta a sobrecarga de informação, e a do engenheiro de software, que enfrenta os desafios de construir sistemas complexos, manuteníveis e escaláveis.

A solução foi materializada como um monólito modular, orientado a eventos, que implementa um pipeline de processamento desacoplado. A abstração do conceito de "vídeo" para uma entidade de "Mídia" genérica, combinada com a comunicação assíncrona via eventos, permitiu a criação de um sistema flexível, capaz de processar diferentes formatos de entrada (vídeo, áudio, texto) através de processadores especializados e independentes.

Os objetivos específicos, que representavam os pilares arquiteturais do projeto, foram alcançados com sucesso. A aplicação rigorosa do padrão \textbf{CQRS} em todos os módulos de negócio resultou em um código mais limpo e focado. O uso do \textbf{Padrão State} garantiu a integridade do ciclo de vida da entidade \texttt{Media}, enquanto os padrões \textbf{Strategy} e \textbf{Factory} permitiram a criação de um sistema de provedores dinâmicos, conferindo alta flexibilidade tecnológica à plataforma. A resiliência foi endereçada com o \textbf{Circuit Breaker Pattern}, e a observabilidade completa foi alcançada através da integração de logs estruturados, métricas Prometheus e tracing distribuído com OpenTelemetry, o que se provou essencial para a depuração de um sistema assíncrono.

A principal contribuição deste trabalho reside na síntese e aplicação prática de um conjunto de padrões de design e arquitetura de software modernos para resolver um problema do mundo real. O projeto demonstra que, ao priorizar o baixo acoplamento, a alta coesão e a separação de responsabilidades desde o início, é possível construir sistemas complexos que são, ao mesmo tempo, robustos, escaláveis e, crucialmente, fáceis de manter e evoluir.

Como trabalhos futuros, vislumbra-se a expansão do pipeline com novos processadores, como um serviço de análise de sentimento sobre o texto transcrito ou um serviço de reconhecimento de objetos em vídeos. A arquitetura atual, por ser modular e orientada a eventos, está perfeitamente preparada para acomodar essas novas funcionalidades com o mínimo de atrito. Além disso, a estratégia de testes pode ser expandida para incluir testes de carga e performance, validando a escalabilidade do sistema sob alta demanda. Finalmente, a infraestrutura de implantação poderia evoluir de Docker Compose para uma plataforma de orquestração mais robusta, como Kubernetes, para a qual a arquitetura baseada em contêineres já está preparada.
 % Capítulo de Conclusão (Novo)

\section*{Agradecimentos}
    Agradeço ao meu orientador Jowaner Araújo pelo apoio técnico e acompanhamento constante durante o desenvolvimento do projeto. Também agradeço à minha família e amigos pelo incentivo emocional ao longo do curso e do TCC.
 % Agradecimentos


% --- Pós-Texto ---
\bibliography{referencias}

\end{document}
